\section{Bargaining anticipates who inherits}
\label{sec:inheritance}

\subsection{Third parties matter to bilateral trades}

With intact creditor inheritance, keeping another player solvent can preserve a blocker against a third player. Conversely, giving a near-bankrupt opponent cash can postpone an estate the buyer expects to receive. These are different strategic reasons for a similar transaction. The following incidents show agents expressing both, with later actions that clarify what they actually do.

\TranscriptTitle{g13}{An estate-insurance pitch selectively presents financing}
At G13 turn 107, Grok has \$35, owns yellow and stands nine squares before Sol's Short Line, where rent would be \$200. Grok offers mortgaged Pennsylvania to Opus for \$200:
\TranscriptQuote{g13-opening-insurance}
Opus counters at \$130, expressing a desire to keep Grok solvent after receiving that appeal. Grok counters at \$200. Its applied note recognizes an available \$130 Ventnor mortgage; one returned provider explanation describes omitting it:
\TranscriptQuote{g13-omission-explanation}
Opus accepts and echoes the stakes:
\TranscriptQuote{g13-acceptance}
\TranscriptFinding{Opus pays \$200 plus \$16 transfer interest, then completes green through an auction seven turns later. In the retained conditional acceptance check, \$130 sale proceeds, \$35 cash and the available Ventnor mortgage would cover the threatened rent. Grok does not land on Short Line on its next roll; it later goes bankrupt there at turn 150, transferring yellow to Sol.}
\TranscriptReading{The seller explicitly frames the deed as protection against a rival's inheritance and describes selective disclosure of collateral. The buyer adopts the estate framing and pays the asking price, without affirming the claimed financing minimum. The text supports omission and uptake of the strategic stakes, not identification of why the extra \$70 was paid. The later estate also shows why a payment does not establish how much survival it purchased.}
\TranscriptSource{G13 d250--253 and d311. Opus uses recovery; its explanations are provider-returned text. Messages are in the actual recipient requests, alongside a faulty past-trade line and correct current holdings. The opening roll claim is overstated for a single roll; no full-turn landing probability is estimated.}

\TranscriptTitle{g18}{An assurance is withdrawn, then bidding resumes lower}
Opus repeatedly offers \$250 for Sol's mortgaged Pacific, which would complete its greens. It renews the assurance after Sol pays it \$750 rent:
\TranscriptQuote{g18-assurance}
Another \$750 payment forces Sol to sell buildings and mortgage deeds, leaving \$49. Sol invokes the price:
\TranscriptQuote{g18-invocation}
Opus withdraws:
\TranscriptQuote{g18-withdrawal}
Its returned explanation says that the payment would prolong a nearly bankrupt rival and that green development would cost too much.
\TranscriptFinding{No sale occurs at \$250. Opus redeems its two green deeds within four turns, then bids \$150 and \$225 for Pacific. Sol asks \$450, \$350 and later \$300; none settles. Opus eventually inherits Sol's estate and wins. The earlier assurance remains in the delivered context when invoked and withdrawn.}
\TranscriptReading{The buyer changes terms after the seller's financial position deteriorates, while continuing to pursue the same group. This is documented withdrawal of a repeated public assurance and subsequent bargaining below that price. It is not proof that the original promise was knowingly false, and no costly reliance is identified. The expressed concern about prolonging a rival connects the negotiation to the creditor's alternative route to acquisition.}
\TranscriptSource{G18 d421--461. Opus uses action-only recovery; the explanation is provider-returned text.}

\subsection{The economic continuation}

Estates are a substantial route of asset transfer in the recorded games. Trade moves 207 deeds; nonterminal bankruptcy moves 227, followed by 170 at the final eliminations. Nonterminal estates complete fifty-two color groups, compared with sixty-three through trade. Sixty-nine further completions occur at game end and have no subsequent earning opportunity. Fourteen of the fifty-two nonterminal estate-completed groups receive rent. This separation prevents terminal bookkeeping from being mistaken for productive acquisition.

The eventual winner receives the first estate in eighteen games, thirty-five of forty nonterminal estates and 202 of their 227 deeds. These outcome-conditioned counts describe a concentration path. They do not show that inheritance independently caused victory. A new check of the delivered state immediately before the first elimination finds its creditor already strictly ahead in building levels in twelve games, cash in eight and deed count in six. Neither ``the strongest player simply inherited'' nor ``inheritance created the strongest player'' is established by these summaries. The initial position and ensuing transfer coexist in the observed trajectory.

The connection to development is direct: thirty-eight of forty nonterminal eliminations follow rent on the creditor's built color property, as do fifty-seven of all sixty eliminations. Such rent can produce cash, an estate, or both residual cash and an estate. Consequently, a receipts-only return calculation misses an economically consequential endpoint. Conversely, an estate with no later earning opportunity should not be counted as realized rent-generating success.

G15 illustrates this sequence without trade by the eventual winner. Gemini completes orange by buying Tennessee at turn 47, builds immediately, and later receives \$3,250 developed orange rent. Orange triggers the first two eliminations; inherited brown triggers the last. Green completes only in the terminal estate. Sol likewise wins G14 without settling a trade. These trajectories show why negotiation volume is not a necessary path to victory in this ruleset, and why the identity of the eventual owner can differ from the player who assembled or first developed a group.
